\documentclass[11pt,a4paper]{article}
\usepackage[utf8]{inputenc}
\usepackage[T1]{fontenc}
\usepackage[spanish,es-tabla]{babel}
\usepackage{mathptmx}
\usepackage[margin=2.3cm]{geometry}
\usepackage{booktabs,tabularx,array}
\usepackage{hyperref}
\newcolumntype{L}{>{\raggedright\arraybackslash}X}

\title{\textbf{Comprensión del contexto y delimitación de la audiencia técnica}\\[2pt]
\large Proyecto 2 --- Visualización I (versión corregida)}
\author{Arnold Santiago Torres Anzola \and Juan David Cruz Romero \and Daniel Alejandro Cipagauta Palacios}
\date{}

\begin{document}
\maketitle

\section{Comprensión del contexto}

El \textit{World Happiness Report 2024} (WHR 2024) compara el bienestar subjetivo de los países a partir de la
evaluación de vida recogida por la Encuesta Mundial de Gallup. Cada persona ubica su vida actual en la escalera
de Cantril, de 0 (peor vida posible) a 10 (mejor vida posible). El informe ordena 143 países según el promedio
nacional de esta evaluación para el trienio 2021--2023.

Para describir esas diferencias, el WHR estima un modelo de regresión con seis indicadores: PIB per cápita
(logaritmo), apoyo social, esperanza de vida saludable, libertad para tomar decisiones, generosidad y percepción
de corrupción. El archivo oficial \textit{Data for Figure 2.1} no contiene los valores originales de esos
indicadores, sino la \textbf{porción de la evaluación de vida de cada país que el modelo asocia con cada
indicador}, medida en puntos de la escalera y en relación con \textit{Distopía}, un país hipotético con los
valores más bajos observados. Estas contribuciones describen asociaciones estadísticas, no efectos causales.

Aunque el WHR 2024 dedica un capítulo al bienestar por grupos de edad, esas desagregaciones \textbf{no forman
parte} del archivo obligatorio. Por eso, este proyecto no compara cohortes ni generaciones: compara países.

\section{Audiencia técnica}

\textbf{Destinatario:} la unidad de análisis comparativo del bienestar de un organismo internacional, es decir,
el equipo de investigadores y analistas de datos encargado de elaborar los diagnósticos de país que acompañan
sus informes de bienestar y desarrollo.

\textbf{Interés en la problemática:} esta unidad necesita saber \emph{en qué dimensiones} se separan los países
con mayor y menor evaluación de vida. Si la brecha se concentrara solo en el PIB per cápita, bastaría con
seguir los indicadores económicos. Si también aparece en dimensiones sociales e institucionales (apoyo social,
libertad para tomar decisiones y percepción de corrupción), el equipo debe incluirlas en sus diagnósticos de país.

\textbf{Qué espera del informe:} definiciones explícitas de cada variable y de su escala, criterios
reproducibles para agrupar países, tratamiento documentado de los valores faltantes y gráficos que permitan
contrastar grupos y localizar países atípicos que merezcan un estudio de caso.

\section{Planteamiento del problema y pregunta analítica}

Los países difieren mucho en la evaluación que sus habitantes hacen de su vida: en el WHR 2024 el promedio va
de 1,7 (Afganistán) a 7,7 (Finlandia). La unidad de análisis necesita saber si esa brecha se concentra en la
dimensión económica o si también aparece en dimensiones de salud y socioinstitucionales, porque de eso depende
qué indicadores incluye en sus diagnósticos de país. El WHR 2024 descompone la evaluación de vida de cada país
en contribuciones asociadas a seis indicadores, lo que permite comparar esas dimensiones en una misma escala
(puntos de la escalera de Cantril).

\begin{quote}
\textbf{\textit{¿En qué dimensiones del bienestar ---económica, de salud y socioinstitucional--- se observan las
mayores diferencias entre los países con mayor y menor evaluación de vida en el World Happiness Report 2024?}}
\end{quote}

\textbf{Subpregunta:} ¿qué países se apartan del patrón de su grupo?

\textbf{Dimensiones:} económica (PIB per cápita), de salud (esperanza de vida saludable) y socioinstitucional
(apoyo social, libertad para tomar decisiones y percepción de corrupción).

\textbf{Justificación de los grupos:} los países ``con mayor y menor evaluación de vida'' son el cuartil
superior (Q4) y el cuartil inferior (Q1) de la variable \texttt{Ladder score} del archivo oficial (cerca de 36
países por grupo). Así se evitan clasificaciones externas que el archivo no contiene, como regiones o niveles
de ingreso, y cualquier persona puede reproducir la agrupación.

\textbf{Alcance:} la pregunta es descriptiva y comparativa. Las diferencias en las contribuciones se
interpretarán como \emph{asociaciones} con la evaluación de vida, no como efectos causales. La generosidad se
reportará como referencia, porque su contribución es pequeña, y \textit{Distopía + residual} porque concentra
lo que el modelo no asocia con ningún indicador.

\section{Variables del archivo oficial (WHR24\_Data\_Figure\_2.1.xls)}

\begin{table}[h]
\centering
\caption{Variables disponibles en la fuente obligatoria y su uso en el análisis.}
\small
\begin{tabularx}{\textwidth}{@{}p{3.6cm}LL@{}}
\toprule
\textbf{Variable (columna)} & \textbf{Qué contiene y escala} & \textbf{Uso en el análisis} \\
\midrule
\texttt{Country name} & Nombre del país (143 países). & Identificador; etiquetas en gráficos interactivos. \\
\texttt{Ladder score} & Promedio nacional de la escalera de Cantril (0--10), trienio 2021--2023. & Variable de referencia; base para construir los cuartiles. \\
\texttt{upperwhisker}, \texttt{lowerwhisker} & Límites del intervalo de confianza al 95\,\% del \texttt{Ladder score}. & Indicar la incertidumbre al comparar países cercanos. \\
\texttt{Explained by: Log GDP per capita} & Puntos de la escalera que el modelo asocia con el PIB per cápita (PPA, en logaritmo), en relación con Distopía. & Dimensión económica. \\
\texttt{Explained by: Healthy life expectancy} & Puntos asociados con la esperanza de vida saludable al nacer. & Dimensión de salud. \\
\texttt{Explained by: Social support} & Puntos asociados con la proporción de personas que dicen tener a quién acudir ante un problema. & Dimensión socioinstitucional. \\
\texttt{Explained by: Freedom to make life choices} & Puntos asociados con la proporción de personas satisfechas con su libertad para decidir sobre su vida. & Dimensión socioinstitucional. \\
\texttt{Explained by: Perceptions of corruption} & Puntos asociados con la percepción de corrupción en el gobierno y en las empresas. Un valor mayor indica \emph{menos} corrupción percibida respecto de Distopía. & Dimensión socioinstitucional. \\
\texttt{Explained by: Generosity} & Puntos asociados con la donación a causas benéficas, ajustada por el PIB per cápita. & Referencia (contribución pequeña). \\
\texttt{Dystopia + residual} & Valor de Distopía más la parte de la evaluación que el modelo no asocia con ningún indicador. & Referencia; señala lo que el modelo no captura. \\
\bottomrule
\end{tabularx}
\end{table}

\noindent\textbf{Notas de preparación.}
(1) El archivo tiene 143 países y 11 columnas. En cada país, la suma de las seis contribuciones y de
\texttt{Dystopia + residual} equivale al \texttt{Ladder score} (diferencia máxima de 0,00005, atribuible al
redondeo); el Notebook verifica esta identidad antes de cualquier análisis.
(2) Bahrain, Tajikistan y State of Palestine tienen la evaluación de vida pero no las contribuciones. Se
mantuvieron en la construcción de los cuartiles y se excluyeron solo de los análisis por dimensión (140 países).
(3) Venezuela registra una contribución del PIB igual a 0 y el Líbano un valor negativo de
\texttt{Dystopia + residual}; ambos casos se señalan al interpretar los resultados.
(4) Las contribuciones están en puntos de la escalera, por lo que se pueden comparar entre dimensiones sin
estandarizarlas. La dimensión socioinstitucional se calculó como la suma de sus tres contribuciones.

\section{Principales hallazgos}

Entre los países del cuartil superior (Q4, media 6,88) y los del cuartil inferior (Q1, media 3,91) hay una
diferencia de \textbf{2,96 puntos} en la evaluación de vida. La Tabla~\ref{tab:brechas} muestra cómo se reparte
esa diferencia entre los componentes; como las contribuciones suman la evaluación de vida, los porcentajes
suman 100\,\%.

\begin{table}[h]
\centering
\caption{Contribución media por cuartil y brecha Q4 $-$ Q1 (puntos de la escalera; 140 países).}
\label{tab:brechas}
\small
\begin{tabular}{@{}llrrrr@{}}
\toprule
\textbf{Dimensión} & \textbf{Componente} & \textbf{Q1} & \textbf{Q4} & \textbf{Brecha} & \textbf{\% de la brecha} \\
\midrule
Económica & PIB per cápita & 0,92 & 1,82 & 0,90 & 30\,\% \\
Socioinstitucional & Apoyo social & 0,74 & 1,41 & 0,67 & 22\,\% \\
 & Libertad para decidir & 0,49 & 0,74 & 0,25 & 8\,\% \\
 & Percepción de corrupción & 0,12 & 0,29 & 0,17 & 6\,\% \\
 & \textit{Subtotal socioinstitucional} & \textit{1,35} & \textit{2,44} & \textit{1,09} & \textit{37\,\%} \\
Salud & Esperanza de vida saludable & 0,35 & 0,66 & 0,32 & 11\,\% \\
Referencia & Generosidad & 0,15 & 0,17 & 0,02 & 1\,\% \\
 & Distopía + residual & 1,15 & 1,79 & 0,65 & 22\,\% \\
\midrule
\multicolumn{2}{@{}l}{\textbf{Evaluación de vida}} & \textbf{3,91} & \textbf{6,88} & \textbf{2,96} & \textbf{100\,\%} \\
\bottomrule
\end{tabular}
\end{table}

\begin{enumerate}
\item \textbf{El PIB per cápita es el componente individual con mayor brecha} (0,90 puntos; 30\,\%), seguido
  por el \textbf{apoyo social} (0,67; 22\,\%). Ambos aumentan de forma gradual de Q1 a Q4 y, entre los extremos,
  sus distribuciones casi no se superponen.
\item \textbf{La dimensión socioinstitucional en conjunto supera a la económica}: apoyo social, libertad y
  corrupción suman 1,09 puntos (37\,\%). La salud aporta 0,32 (11\,\%) y la generosidad no diferencia a los
  grupos: sus intervalos de confianza al 95\,\% se superponen, mientras que en los otros cinco componentes no.
\item \textbf{La percepción de corrupción solo distingue al grupo más alto}: es similar en Q1, Q2 y Q3
  (alrededor de 0,1) y sube únicamente en Q4, donde además varía mucho entre países (de 0 a 0,57).
\item \textbf{Las dimensiones no son independientes.} PIB, salud y apoyo social varían juntos (correlaciones de
  Spearman entre 0,72 y 0,86), por lo que sus brechas no pueden leerse como efectos separados. Libertad y
  corrupción se relacionan menos con el PIB (0,44 y 0,38) y aportan información que el PIB no resume.
\item \textbf{Hay países que se apartan del patrón de su grupo.} Con una contribución del PIB propia de Q1,
  Kirguistán (Q2), Uzbekistán (Q3) y Nicaragua (Q3) tienen niveles socioinstitucionales altos. En sentido
  contrario, el Líbano (Q1) y Turquía (Q2) tienen contribuciones del PIB propias de Q3 o Q4, pero niveles
  socioinstitucionales bajos.
\item \textbf{Una parte importante de la brecha no se asocia con ningún indicador} (Distopía + residual, 0,65
  puntos; 22\,\%). En varios países de Q1, como Níger, Burkina Faso o Chad, es uno de los componentes más grandes.
\end{enumerate}

\section{Conclusiones y limitaciones}

\textbf{Respuesta a la pregunta.} Las mayores diferencias entre los países con mayor y menor evaluación de vida
se observan en la \textbf{dimensión económica (PIB)} y en el \textbf{apoyo social}; consideradas en conjunto, las
dimensiones \textbf{socioinstitucionales} concentran una parte de la brecha mayor que la económica. La salud
muestra una brecha menor y la generosidad no diferencia a los grupos.

\textbf{Implicación para la audiencia.} Un diagnóstico de país centrado solo en indicadores económicos dejaría
fuera una parte de la brecha de bienestar de magnitud comparable o mayor. Apoyo social, libertad y percepción de
corrupción merecen seguimiento propio, en especial libertad y corrupción, que son las menos redundantes con el
PIB. Los países atípicos identificados son candidatos a estudios de caso.

\textbf{Limitaciones.} (1) Las contribuciones son asociaciones estimadas por el modelo del WHR, no efectos
causales. (2) Cada contribución es el indicador del país multiplicado por un coeficiente común, de modo que la
comparación entre dimensiones refleja también esos pesos. (3) Los datos son promedios nacionales de 2021--2023:
no muestran desigualdades dentro de los países ni cambios en el tiempo. (4) Los cuartiles son una elección del
equipo y Q1 es muy heterogéneo (de 1,72 a 4,66 puntos).

\textbf{Nuevas preguntas.} ¿Se mantiene el patrón con los valores originales de los indicadores? ¿Qué
características comparten los países que se apartan de su grupo? ¿Cómo cambian estas brechas por grupos de edad
o a lo largo del tiempo? (Las dos últimas requieren datos que no están en el archivo de la Figura 2.1.)

\section*{Referencias}
\small
Helliwell, J. F., Layard, R., Sachs, J. D., De Neve, J.-E., Aknin, L. B. y Wang, S. (Eds.). (2024).
\textit{World Happiness Report 2024}. Wellbeing Research Centre, University of Oxford.
\url{https://www.worldhappiness.report/ed/2024/}

\noindent Datos: \textit{Data for Figure 2.1}, WHR 2024.
\url{https://files.worldhappiness.report/WHR24_Data_Figure_2.1.xls}

\end{document}
